% =====================================================================
\chapter{Modulo 14 -- Grafi}
% =====================================================================

\tikzset{
  vert/.style={circle,draw,thick,fill=heapcol,font=\small\bfseries,minimum size=7mm,inner sep=1pt},
  arco/.style={thick},
  arcoO/.style={-{Stealth[length=2.4mm]},thick},
  albArco/.style={line width=2pt,green!50!black},
  nodoG/.style={draw,rectangle split,rectangle split horizontal,rectangle split parts=2,
                fill=heapcol,font=\ttfamily\small,minimum height=0.6cm,
                rectangle split part fill={heapcol,heapcol!60!sfondo}},
  grigio/.style={fill=black!25!sfondo},
  nero/.style={fill=black!65!sfondo,text=sfondo},
}

% ---------------------------------------------------------------------
\section{Teoria \& Analogia Chiave}
% ---------------------------------------------------------------------

\subsection*{Dagli alberi ai grafi}

Ricapitoliamo il percorso delle strutture dati. In una \textbf{lista} ogni nodo ha un
solo successivo: una fila indiana. In un \textbf{albero} ogni nodo può avere più figli,
ma \textbf{un solo padre}, e non si torna mai indietro: una gerarchia. Ora togliamo
anche quest'ultimo vincolo: in un \textbf{grafo} ogni nodo può essere collegato a
\textbf{qualunque} altro nodo, quanti ne vuole, anche formando dei giri chiusi.

\begin{analogia}[: la cartina stradale]
Pensa alla cartina della Sicilia. Le \textbf{città} sono i nodi; le \textbf{strade} che
le collegano direttamente sono gli archi. Da Catania si va a Siracusa, a Messina, a
Enna; da Enna a Palermo; da Palermo a Messina\ldots{} Non c'è una ``radice'', nessuna
città è ``padre'' di un'altra, e partendo da Catania si può tornare a Catania facendo un
giro (Catania, Enna, Palermo, Messina, Catania). Un grafo è esattamente questo: un
insieme di luoghi e di collegamenti diretti tra coppie di luoghi.
\end{analogia}

Altri grafi che usi ogni giorno senza saperlo: la rete di amicizie di un social (persone
e amicizie), la metropolitana (stazioni e tratte), Internet (computer e cavi), le pagine
web (pagine e link), i prerequisiti degli esami (esami e propedeuticità). Per questo le
slide dicono che i grafi sono strutture ``di fondamentale importanza'', con centinaia di
problemi legati a loro. Al corso ne vediamo i più elementari: \textbf{come si memorizza}
un grafo e \textbf{come lo si visita}.

\subsection*{Il vocabolario}

\begin{definizione}[: grafo]
Un \textbf{grafo} $G = (V, E)$ è formato da:
\begin{itemize}
  \item un insieme $V$ di \textbf{vertici} (o nodi);
  \item un insieme $E$ di \textbf{archi}: ogni arco è una coppia di vertici $(u, v)$.
\end{itemize}
Se $(u, v)$ è un arco, si dice che $v$ è \textbf{adiacente} a $u$.
\end{definizione}

La lettera $V$ viene da \emph{vertex}, la $E$ da \emph{edge} (arco, in inglese). Con
$|V|$ si indica il \textbf{numero} di vertici e con $|E|$ il numero di archi: le
barrette vogliono solo dire ``quanti sono''. Spesso si scrive più semplicemente $n$ per
i vertici e $m$ per gli archi.

Ci sono due famiglie di grafi, e la differenza è importante per tutto il resto:
\begin{itemize}
  \item \textbf{Grafo orientato} (o \emph{direzionato}): gli archi hanno un verso, come
        le strade a senso unico. L'arco $(u, v)$ va \emph{da} $u$ \emph{a} $v$: è
        \textbf{uscente} da $u$ ed \textbf{entrante} in $v$. Da $u$ si arriva a $v$, ma
        non è detto il contrario. Si disegna con una freccia. Esempio: ``A segue B'' su un
        social, ``l'esame X è propedeutico a Y''.
  \item \textbf{Grafo non orientato}: gli archi non hanno verso, come le strade a doppio
        senso. $(u, v)$ e $(v, u)$ sono lo stesso arco, e l'adiacenza è
        \textbf{simmetrica}: se $v$ è adiacente a $u$, anche $u$ è adiacente a $v$. Si
        disegna con una linea semplice. Esempio: ``A è amico di B'' (l'amicizia è
        reciproca). Nelle slide i grafi non orientati non hanno \textbf{cappi}
        (\emph{self-loop}), cioè archi da un vertice a sé stesso.
\end{itemize}

Altre parole che ricorrono nei quiz:
\begin{itemize}
  \item \textbf{Grado} di un vertice: quanti archi lo toccano. In un grafo non orientato
        è il numero dei suoi vicini. In un grafo orientato si distinguono \textbf{grado
        entrante} (frecce che arrivano) e \textbf{grado uscente} (frecce che partono), e
        il grado è la loro somma.
  \item \textbf{Cammino} da $u$ a $v$: una sequenza di vertici $v_0, v_1, \ldots, v_k$
        con $v_0 = u$, $v_k = v$, in cui ogni coppia di vertici consecutivi è un arco. La
        sua \textbf{lunghezza} è il numero di archi, $k$. Se esiste un cammino da $u$ a
        $v$, si dice che $v$ è \textbf{raggiungibile} da $u$. Il cammino è
        \textbf{semplice} se non passa due volte dallo stesso vertice.
  \item \textbf{Ciclo}: un cammino che torna al punto di partenza ($v_0 = v_k$). Un
        grafo senza cicli è \textbf{aciclico}. Un grafo orientato aciclico si chiama
        \textbf{DAG} (\emph{Directed Acyclic Graph}).
  \item Un grafo non orientato è \textbf{connesso} se ogni coppia di vertici è unita da un
        cammino: ``da ogni città si arriva a ogni altra''. Se non lo è, si spezza in
        \textbf{componenti connesse}: i gruppi di vertici che si raggiungono tra loro
        (pensa alle isole: la Sicilia e le Eolie sono due componenti, se contano solo le
        strade).
  \item In un grafo orientato si parla di \textbf{fortemente connesso} se per ogni coppia
        $u, v$ c'è un cammino da $u$ a $v$ \emph{e} uno da $v$ a $u$.
  \item Un grafo non orientato è \textbf{completo} se ogni coppia di vertici distinti è
        adiacente: con $n$ vertici ha $n(n-1)/2$ archi (Esercizio~14.1).
  \item $G' = (V', E')$ è un \textbf{sottografo} di $G$ se prende solo alcuni dei vertici e
        alcuni degli archi di $G$.
\end{itemize}

\textbf{Un albero è un grafo particolare}: non orientato, connesso e senza cicli. Tutto
quello che diremo sui grafi vale anche per gli alberi; il viceversa no.

\subsection*{Come si memorizza un grafo: due modi}

Dentro il computer i vertici sono di solito \textbf{numerati} da $0$ a $n - 1$: il
vertice ``Catania'' diventa, ad esempio, il numero 2. Si tiene un array di chiavi
(\cc|chiavi[2] = "Catania"|) e da lì in poi si lavora con gli indici, che sono comodi
per indicizzare array. Restano da memorizzare gli archi, e qui le strade sono due.

\paragraph{1. Matrice di adiacenza.}
Una tabella $n \times n$ di \cc|bool|: la cella \cc|adj[i][j]| vale \cc|true| (1) se c'è
l'arco da $i$ a $j$, \cc|false| (0) altrimenti.

\begin{analogia}[: la tabella delle distanze dell'atlante stradale]
Negli atlanti stradali c'è una tabella con le città sia sulle righe sia sulle colonne:
all'incrocio tra la riga ``Catania'' e la colonna ``Palermo'' leggi la distanza. La
matrice di adiacenza è la stessa tabella, ma invece della distanza c'è solo un sì o un
no: ``c'è una strada diretta?''. Per rispondere basta andare all'incrocio giusto, senza
cercare: un'operazione istantanea, $O(1)$.
\end{analogia}

Proprietà:
\begin{itemize}
  \item in un grafo \textbf{non orientato} la matrice è \textbf{simmetrica} rispetto alla
        diagonale: ogni arco si scrive due volte, in \cc|adj[i][j]| e in \cc|adj[j][i]|;
  \item la somma della riga $i$ è il grado del vertice $i$;
  \item occupa sempre $n \times n$ celle, anche se gli archi sono pochissimi: spazio
        $O(|V|^2)$. Con 10\,000 vertici sono 100 milioni di celle, quasi tutte zero;
  \item per trovare i vicini di $u$ bisogna scorrere tutta la riga: $n$ passi, anche se
        i vicini sono due.
\end{itemize}

In C++ è un array dinamico di array dinamici, come le matrici del Modulo 3. È la
rappresentazione usata dal codice del professore (\cc|bool** adj|):
\begin{lstlisting}
bool** adj = new bool*[n];             // un array di n puntatori a riga
for (int i = 0; i < n; i++)
    adj[i] = new bool[n]{};            // ogni riga: n celle a false
adj[i][j] = true;                      // aggiungo l'arco (i, j)
adj[j][i] = true;                      // ...e (j, i) se il grafo NON è orientato
\end{lstlisting}

\paragraph{2. Liste di adiacenza.}
Un array di $n$ \textbf{liste concatenate}, una per vertice: la lista \cc|adj[u]| contiene
i vertici adiacenti a $u$ (i suoi vicini). Ogni elemento è un nodo dinamico come quelli
del Modulo 11:
\begin{lstlisting}
struct Arco {
    int dest;          // indice del vertice di arrivo
    Arco* succ;        // prossimo vicino nella lista
};
Arco** adj = new Arco*[n];             // n teste di lista
for (int i = 0; i < n; i++) adj[i] = nullptr;      // all'inizio nessun vicino
adj[u] = new Arco{v, adj[u]};          // arco (u, v): inserimento in testa, O(1)
adj[v] = new Arco{u, adj[v]};          // ...e (v, u) se il grafo NON è orientato
\end{lstlisting}

\begin{analogia}[: la rubrica degli amici]
Ognuno ha la sua rubrica con i soli numeri dei propri amici. Per sapere chi sono gli
amici di Anna apri la rubrica di Anna e la leggi tutta: è lunga quanto il numero dei suoi
amici, non quanto il numero di abitanti della città. Ma per sapere se Anna conosce Bruno
devi sfogliarla tutta, perché non c'è una casella ``Bruno'' già pronta.
\end{analogia}

Proprietà (dalle slide):
\begin{itemize}
  \item in un grafo \textbf{non orientato} ogni arco compare in \textbf{due} liste, quindi
        la somma delle lunghezze di tutte le liste è $2|E|$; in un grafo
        \textbf{orientato} ogni arco compare una volta sola, e la somma è $|E|$;
  \item spazio: l'array di $|V|$ teste più i nodi delle liste, cioè $O(|V| + |E|)$. È
        molto meglio della matrice per i grafi \textbf{sparsi} (con pochi archi), che nella
        pratica sono la maggioranza;
  \item scorrere i vicini di $u$ costa quanto il numero dei vicini: il meglio possibile;
  \item sapere se $(u, v)$ è un arco costa la lunghezza della lista di $u$.
\end{itemize}

\begin{center}
\begin{tabular}{lcc}
\toprule
\textbf{Operazione} & \textbf{Matrice} & \textbf{Liste} \\
\midrule
spazio occupato                    & $O(|V|^2)$     & $O(|V| + |E|)$ \\
\q{c'è l'arco $(u,v)$?}            & $O(1)$         & $O(\text{grado di } u)$ \\
scorrere i vicini di $u$           & $O(|V|)$       & $O(\text{grado di } u)$ \\
aggiungere un arco                 & $O(1)$         & $O(1)$ (in testa) \\
visita BFS o DFS completa          & $O(|V|^2)$     & $O(|V| + |E|)$ \\
\bottomrule
\end{tabular}
\end{center}
\textbf{Regola pratica}: grafo piccolo o ``denso'' (quasi tutti collegati con tutti) e
tante domande ``$u$ e $v$ sono collegati?'' $\Rightarrow$ matrice. Grafo grande e
sparso, e algoritmi che scorrono i vicini (le visite) $\Rightarrow$ liste. Le slide
assumono le liste per BFS e DFS.

\subsection*{Visitare un grafo: perché serve ``ricordarsi'' dove si è già stati}

Visitare un grafo significa, come per gli alberi, raggiungere tutti i vertici partendo da
uno. Ma c'è una differenza enorme: \textbf{nei grafi ci sono i cicli}. Se da A vado a B,
da B a C e da C di nuovo ad A, senza un promemoria continuerei a girare in tondo per
sempre. Negli alberi questo problema non esisteva, perché scendendo non si torna mai
indietro.

La soluzione delle slide è \textbf{colorare} i vertici:
\begin{itemize}
  \item \textbf{bianco}: non ancora scoperto (non l'ho mai visto);
  \item \textbf{grigio}: scoperto, ma non ho ancora finito di esaminare i suoi vicini (è
        la ``frontiera'' tra il noto e l'ignoto);
  \item \textbf{nero}: scoperto e completamente esaminato.
\end{itemize}
Un vertice si visita solo se è \textbf{bianco}. Così ogni vertice viene scoperto
\textbf{una volta sola}, e la visita termina sempre. Per molti esercizi bastano due colori
(un \cc|bool visitato[]|); il grigio serve in casi come la ricerca di cicli
(Esercizio~14.8).

\begin{analogia}[: il gesso nel labirinto]
Chi esplora una grotta segna col gesso ogni incrocio in cui passa. Arrivato a un incrocio
già segnato, sa di esserci già stato e non lo riesplora. Il colore dei vertici è il gesso.
\end{analogia}

Esistono due strategie di visita, che differiscono nell'\textbf{ordine} in cui si
esplora: in \textbf{ampiezza} (BFS) e in \textbf{profondità} (DFS).

\subsection*{BFS: la visita in ampiezza (\emph{Breadth-First Search})}

\begin{analogia}[: il sasso nello stagno, o il passaparola]
Lanci un sasso in uno stagno: le onde si allargano in cerchi. Prima viene raggiunto
tutto ciò che sta a distanza 1, poi tutto ciò che sta a distanza 2, e così via. Oppure
pensa a una notizia: la dici ai tuoi amici (distanza 1), loro la dicono ai loro amici
(distanza 2, gli ``amici di amici''), e così via. La BFS visita il grafo \textbf{a
cerchi concentrici} attorno alla sorgente.
\end{analogia}

Per realizzare i cerchi serve una \textbf{coda} (FIFO, Modulo 12): chi è scoperto prima
viene esaminato prima. È esattamente la visita per livelli degli alberi del Modulo 13, con
in più i colori per non ripassare dagli stessi vertici. Ecco l'algoritmo delle slide,
$\text{BFS}(G, s)$ con sorgente $s$:

\begin{passopasso}[: BFS(G, s)]
\begin{enumerate}
  \item Ogni vertice $u$ diventa \textbf{bianco}, con distanza $d[u] = \infty$ (``non
        so ancora'') e predecessore $\text{pred}[u] = \text{NULL}$.
  \item La sorgente $s$ diventa \textbf{grigia}, con $d[s] = 0$. Metti $s$ nella coda.
  \item Finché la coda non è vuota:
  \begin{enumerate}
    \item estrai dalla testa della coda un vertice $u$;
    \item per ogni vicino $v$ di $u$ (scorri la lista \cc|adj[u]|): se $v$ è
          \textbf{bianco}, l'hai appena scoperto: coloralo di grigio, poni
          $d[v] = d[u] + 1$ (è un passo più lontano di $u$), $\text{pred}[v] = u$ (``ci
          sono arrivato da $u$''), e mettilo in fondo alla coda;
    \item $u$ diventa \textbf{nero}: hai finito con lui.
  \end{enumerate}
\end{enumerate}
\end{passopasso}

Le slide fanno notare l'\textbf{invariante}: in ogni momento la coda contiene esattamente
i vertici grigi. E poiché la coda è FIFO, i vertici escono in ordine di distanza: prima
tutti quelli a distanza 0, poi quelli a distanza 1, poi a distanza 2\ldots

\textbf{Cosa calcola la BFS}, oltre all'ordine di visita:
\begin{itemize}
  \item $d[v]$ è la \textbf{distanza minima} da $s$ a $v$, cioè il minimo numero di archi
        da percorrere. Vale per i grafi \textbf{non pesati} (archi tutti ``uguali''): se le
        strade avessero lunghezze diverse servirebbero altri algoritmi (Dijkstra, fuori
        programma);
  \item i puntatori $\text{pred}[v]$ formano un albero, il \textbf{breadth-first tree}:
        risalendo da $v$ con i pred si torna a $s$ lungo un \textbf{cammino minimo};
  \item i vertici rimasti bianchi (con $d = \infty$) \textbf{non sono raggiungibili} da $s$.
\end{itemize}

Per stampare il cammino minimo da $s$ a $v$ le slide danno una procedura
\textbf{ricorsiva}, \emph{Print-Path}: il cammino fino a $v$ è ``il cammino fino a
$\text{pred}[v]$, seguito da $v$''. Le tre domande del Modulo 8: caso base $v = s$
(stampo $s$); se $\text{pred}[v]$ è NULL non c'è cammino; altrimenti chiamo su
$\text{pred}[v]$ e \textbf{dopo} stampo $v$ (la stampa dopo la chiamata rovescia l'ordine,
come nel Modulo 8). Costo $O(|V|)$: ogni chiamata accorcia il cammino di un vertice.
\begin{lstlisting}
void stampaCammino(int s, int v) {
    if (v == s) cout << chiavi[s];
    else if (pred[v] == -1) cout << "nessun cammino";
    else { stampaCammino(s, pred[v]); cout << " -> " << chiavi[v]; }
}
\end{lstlisting}

\subsection*{DFS: la visita in profondità (\emph{Depth-First Search})}

\begin{analogia}[: il filo di Arianna]
Teseo nel labirinto: a ogni incrocio prende la prima strada non ancora esplorata e va
\textbf{avanti più che può}, srotolando il filo. Quando arriva in un vicolo cieco (o in
un posto già visitato) \textbf{torna indietro} seguendo il filo fino all'ultimo incrocio
che ha ancora strade nuove, e riparte da lì. Questo ``tornare indietro'' si chiama
\textbf{backtracking}.
\end{analogia}

Il ``filo'' è una \textbf{pila}: l'ultimo incrocio raggiunto è il primo a cui si torna
(LIFO). E la pila più comoda che abbiamo è quella delle chiamate ricorsive (Modulo 8):
per questo la DFS si scrive quasi sempre in modo \textbf{ricorsivo}, ed è cortissima.
L'algoritmo delle slide ha due parti:

\begin{passopasso}[: DFS(G) e DFS-Visit(u)]
\textbf{DFS(G)}:
\begin{enumerate}
  \item ogni vertice diventa bianco, con pred NULL; poni $\text{tempo} = 0$;
  \item per ogni vertice $u$, \textbf{se è ancora bianco}, chiama DFS-Visit($u$).
\end{enumerate}
\textbf{DFS-Visit(u)}:
\begin{enumerate}
  \item $u$ diventa grigio; $d[u] = \text{tempo}$\texttt{++} (istante di
        \textbf{scoperta});
  \item per ogni vicino $v$ di $u$: se $v$ è bianco, $\text{pred}[v] = u$ e chiama
        \textbf{ricorsivamente} DFS-Visit($v$);
  \item $u$ diventa nero; $f[u] = \text{tempo}$\texttt{++} (istante di \textbf{fine}).
\end{enumerate}
\end{passopasso}

Tre osservazioni:
\begin{itemize}
  \item il ciclo esterno di DFS($G$) fa ripartire la visita da ogni vertice non ancora
        raggiunto: così la DFS visita \textbf{tutto} il grafo, anche se non è connesso.
        Ogni ripartenza produce un albero diverso: l'insieme dei pred forma una
        \textbf{foresta} (\emph{depth-first forest}). La BFS invece parte da una sola
        sorgente;
  \item i \textbf{timestamp} $d[u]$ e $f[u]$ sono come le timbrature di entrata e uscita
        di un cartellino: $d[u] < f[u]$ sempre, e se $v$ è scoperto dentro la visita di
        $u$ allora il suo intervallo $[d[v], f[v]]$ sta dentro quello di $u$, proprio come
        le chiamate ricorsive annidate;
  \item la DFS \textbf{non} calcola le distanze minime: va lontano subito, e può
        arrivare a un vicino della sorgente passando per una strada lunghissima.
\end{itemize}

Le slide classificano anche gli archi incontrati dalla DFS: \textbf{tree edge} (archi
dell'albero, usati per scoprire un vertice), \textbf{back edge} (verso un antenato ancora
grigio: segnalano un \textbf{ciclo}), \textbf{forward edge} (verso un discendente già
finito), \textbf{cross edge} (tutti gli altri). Per l'esame basta ricordare che un arco
verso un vertice \textbf{grigio} chiude un ciclo (Esercizio~14.8).

\subsection*{Un'applicazione della DFS: l'ordinamento topologico}

Hai un elenco di attività con dei vincoli ``A va fatta prima di B'': gli esami
propedeutici, o l'esempio delle slide dei vestiti (le calze prima delle scarpe, la
camicia prima della cravatta). Le attività sono i vertici, ogni vincolo è un arco
\textbf{orientato} da A a B. Un \textbf{ordinamento topologico} è una disposizione in fila
di tutti i vertici in cui ogni arco va \textbf{da sinistra a destra}: un ordine in cui
fare le attività rispettando tutti i vincoli.

Esiste solo se il grafo è un \textbf{DAG} (orientato e aciclico): se A richiede B e B
richiede A, nessun ordine funziona. Il trucco delle slide:
\begin{passopasso}[: TOPOLOGICAL-SORT(G)]
\begin{enumerate}
  \item esegui DFS($G$);
  \item ogni volta che un vertice \textbf{finisce} (diventa nero), inseriscilo \textbf{in
        testa} a una lista;
  \item alla fine la lista contiene i vertici in ordine decrescente di tempo di fine $f$:
        è un ordinamento topologico. Costo $O(|V| + |E|)$.
\end{enumerate}
\end{passopasso}
Perché funziona? Un vertice finisce solo dopo che sono finiti tutti quelli raggiungibili
da lui, cioè tutto ciò che ``viene dopo''. Inserendolo in testa finisce \emph{davanti} a
tutti loro. Lo vedrai al lavoro nello Step 6 del laboratorio.

Le slide accennano infine alle \textbf{componenti fortemente connesse}, calcolabili con due
DFS (una sul grafo, una sul grafo \textbf{trasposto} $G^T$, che ha tutti gli archi
rovesciati, Esercizio~14.4). Basta sapere che esistono e che costano $O(|V|+|E|)$.

\subsection*{Quanto costano le visite (senza formule)}

Usiamo il contatore di passi del Modulo 6. Con le \textbf{liste di adiacenza}:
\begin{itemize}
  \item ogni vertice entra in coda (o viene visitato) \textbf{una volta sola}, grazie ai
        colori: $|V|$ passi;
  \item quando un vertice viene esaminato, si scorre la sua lista \textbf{una volta sola}:
        in totale si scorrono tutte le liste, cioè $2|E|$ elementi (non orientato) o
        $|E|$ (orientato).
\end{itemize}
Totale: circa $|V| + |E|$ passi, cioè $O(|V| + |E|)$, sia per BFS sia per DFS. In
parole: \textbf{proporzionale alla grandezza del grafo}, il meglio possibile, visto che il
grafo bisogna almeno leggerlo.

Con la \textbf{matrice}, per trovare i vicini di un vertice bisogna scorrere tutta la sua
riga ($|V|$ celle), e lo si fa per ogni vertice: $|V| \times |V|$ passi, $O(|V|^2)$.
\begin{center}
\begin{tabular}{lrr}
\toprule
\textbf{Grafo} & \textbf{Liste} ($\approx V + 2E$) & \textbf{Matrice} ($V \times V$) \\
\midrule
100 vertici, 300 archi           & 700 passi       & 10\,000 passi \\
1\,000 vertici, 3\,000 archi     & 7\,000 passi    & 1\,000\,000 passi \\
100\,000 vertici, 300\,000 archi & 700\,000 passi  & 10\,000\,000\,000 passi \\
\bottomrule
\end{tabular}
\end{center}

\subsection*{BFS e DFS a confronto}
\begin{center}
\begin{tabular}{p{3.2cm}p{4.6cm}p{4.6cm}}
\toprule
 & \textbf{BFS (ampiezza)} & \textbf{DFS (profondità)} \\
\midrule
struttura di supporto & coda (FIFO)                     & pila / ricorsione (LIFO) \\
ordine                & a cerchi, per distanza          & più lontano possibile, poi indietro \\
cammini minimi        & sì (grafi non pesati)           & no \\
parte da              & una sorgente $s$                & tutti i vertici (foresta) \\
usi tipici            & distanze, \q{amici di amici}, labirinti & componenti, cicli, ordinamento topologico \\
costo (liste)         & $O(|V| + |E|)$                  & $O(|V| + |E|)$ \\
\bottomrule
\end{tabular}
\end{center}

\begin{inparole}
\begin{itemize}
  \item Grafo = vertici + archi. Orientato (frecce, senso unico) o non orientato
        (doppio senso, adiacenza simmetrica). Possono esserci cicli.
  \item Matrice: tabella $n \times n$ di sì/no, spazio $V^2$, \q{c'è l'arco?} istantaneo.
  \item Liste: per ogni vertice la lista dei vicini, spazio $V + E$, ideale per le visite.
        Non orientato: ogni arco in due liste (somma $2E$).
  \item Si colorano i vertici (bianco, grigio, nero) per non girare in tondo nei cicli.
  \item BFS: coda, cerchi concentrici, distanze minime e cammini con pred + Print-Path.
  \item DFS: ricorsione, va a fondo e torna indietro, tempi $d$/$f$, ripartenza da ogni
        vertice bianco. Ordinamento topologico = vertici in testa a una lista quando
        finiscono.
  \item Entrambe costano $O(V + E)$ con le liste, $O(V^2)$ con la matrice.
\end{itemize}
\end{inparole}

% ---------------------------------------------------------------------
\section{Rappresentazione Visiva della Memoria (TikZ)}
% ---------------------------------------------------------------------

\textbf{Il grafo di esempio} (quello del pattern), non orientato, con 6 vertici e 7 archi,
accanto a un piccolo grafo orientato.
\begin{center}
\begin{tikzpicture}
  \node[vert] (A) at (0,0) {A};
  \node[vert] (B) at (1.6,1) {B};
  \node[vert] (C) at (1.6,-1) {C};
  \node[vert] (D) at (3.2,1) {D};
  \node[vert] (E) at (3.2,-1) {E};
  \node[vert] (F) at (4.8,0) {F};
  \draw[arco] (A)--(B) (A)--(C) (B)--(D) (C)--(D) (C)--(E) (D)--(F) (E)--(F);
  \node[etichetta] at (2.4,-1.9) {non orientato: A--B è anche B--A};
  \begin{scope}[xshift=7.5cm]
    \node[vert] (1) at (0,1) {1};
    \node[vert] (2) at (2,1) {2};
    \node[vert] (3) at (0,-1) {3};
    \node[vert] (4) at (2,-1) {4};
    \draw[arcoO] (1)--(2);
    \draw[arcoO] (2)--(4);
    \draw[arcoO] (4)--(3);
    \draw[arcoO] (3)--(1);
    \draw[arcoO] (1)--(4);
    \node[etichetta] at (1,-1.9) {orientato: c'è $1 \to 2$ ma non $2 \to 1$};
  \end{scope}
\end{tikzpicture}
\end{center}
Nel grafo orientato il vertice 4 ha grado entrante 2 (da 2 e da 1) e uscente 1 (verso 3);
il cammino $1 \to 2 \to 4 \to 3 \to 1$ è un ciclo.

\needspace{7cm}
\textbf{Lo stesso grafo nelle due rappresentazioni.} A sinistra la matrice (simmetrica:
la parte sopra la diagonale è lo specchio di quella sotto). A destra l'array delle teste
di lista: ogni cella punta alla catena dei vicini. Le liste sono in ordine alfabetico
perché il pattern inserisce in fondo; inserendo in testa sarebbero rovesciate.
\begin{center}
\begin{tikzpicture}
  % matrice
  \foreach \c [count=\j from 0] in {A,B,C,D,E,F} {
    \node[font=\ttfamily\small\bfseries] at (0.55*\j+0.55,0.5) {\c};
    \node[font=\ttfamily\small\bfseries] at (0,-0.55*\j) {\c};
  }
  \foreach \riga [count=\i from 0] in {{0,1,1,0,0,0},{1,0,0,1,0,0},{1,0,0,1,1,0},{0,1,1,0,0,1},{0,0,1,0,0,1},{0,0,0,1,1,0}} {
    \foreach \x [count=\j from 0] in \riga {
      \ifnum\x=1
        \node[draw,minimum size=5.5mm,font=\ttfamily\small,fill=green!25!sfondo] at (0.55*\j+0.55,-0.55*\i) {1};
      \else
        \node[draw,minimum size=5.5mm,font=\ttfamily\small,text=black!45!sfondo] at (0.55*\j+0.55,-0.55*\i) {0};
      \fi
    }
  }
  \node[etichetta] at (1.9,-3.6) {\texttt{bool** adj}: 36 celle};
  % liste
  \begin{scope}[xshift=5.2cm]
    \foreach \v/\vicini [count=\i from 0] in {A/{B,C},B/{A,D},C/{A,D,E},D/{B,C,F},E/{C,F},F/{D,E}} {
      \node[cellina,minimum height=0.55cm,minimum width=0.6cm] (t\i) at (0,-0.6*\i) {};
      \node[font=\ttfamily\small\bfseries,anchor=east] at (-0.35,-0.6*\i) {\v};
      \foreach \w [count=\k from 0] in \vicini {
        \node[nodoG,minimum height=0.45cm] (n\i-\k) at (1.4+1.25*\k,-0.6*\i) {\w \nodepart{two} \phantom{x}};
      }
      \draw[freccia,{Circle[length=1.5mm]}-{Stealth[length=2mm]}] (t\i.center) -- (n\i-0.west);
    }
    \foreach \i/\k in {0/1,1/1,2/1,2/2,3/1,3/2,4/1,5/1} {
      \pgfmathtruncatemacro{\km}{\k-1}
      \draw[freccia,{Circle[length=1.5mm]}-{Stealth[length=2mm]}] (n\i-\km.two split south) ++(0.12,0.22) -- (n\i-\k.west);
    }
    \node[etichetta] at (2.3,-3.6) {\texttt{Arco** adj}: 6 teste + 14 nodi $= 2 \times 7$ archi};
  \end{scope}
\end{tikzpicture}
\end{center}
L'ultimo nodo di ogni lista ha \cc|succ = nullptr| (la parte destra vuota). Con 6 vertici
le due rappresentazioni si equivalgono; con 1000 vertici e 3000 archi la matrice avrebbe
un milione di celle, le liste 6000 nodi.

\needspace{8cm}
\textbf{La BFS da A, passo per passo.} A ogni riga: il vertice estratto dalla testa
della coda, i vicini bianchi scoperti (con la loro distanza) e la coda dopo il passo.
\begin{center}
\begin{tabular}{clll}
\toprule
\textbf{Passo} & \textbf{Estraggo} & \textbf{Scopro (bianchi $\to$ grigi)} & \textbf{Coda dopo} \\
\midrule
0 & --- & A ($d=0$)                & [ A ] \\
1 & A   & B ($d=1$), C ($d=1$)     & [ B, C ] \\
2 & B   & D ($d=2$) \quad (A è già nero) & [ C, D ] \\
3 & C   & E ($d=2$) \quad (A nero, D grigio) & [ D, E ] \\
4 & D   & F ($d=3$)                & [ E, F ] \\
5 & E   & --- \quad (C nero, F grigio) & [ F ] \\
6 & F   & ---                      & [ ] fine \\
\bottomrule
\end{tabular}
\end{center}
Il risultato: i vertici disposti a ``cerchi'' di distanza, e in verde gli archi
$\text{pred}$ che formano il breadth-first tree. I due archi grigi chiari (C--D ed E--F)
non servono a scoprire nessuno: collegano vertici già scoperti.
\begin{center}
\begin{tikzpicture}
  \foreach \x/\t in {0/{$d=0$},2/{$d=1$},4/{$d=2$},6/{$d=3$}}
    \node[etichetta] at (\x,1.8) {\t};
  \foreach \x in {1,3,5} \draw[dashed,black!30!sfondo] (\x,1.5) -- (\x,-1.5);
  \node[vert,nero] (A) at (0,0) {A};
  \node[vert,nero] (B) at (2,0.8) {B};
  \node[vert,nero] (C) at (2,-0.8) {C};
  \node[vert,nero] (D) at (4,0.8) {D};
  \node[vert,nero] (E) at (4,-0.8) {E};
  \node[vert,nero] (F) at (6,0) {F};
  \draw[arco,black!25!sfondo] (C)--(D) (E)--(F);
  \draw[albArco] (A)--(B) (A)--(C) (B)--(D) (C)--(E) (D)--(F);
  \node[etichetta,anchor=west,text width=5cm] at (7,0) {Cammino minimo A$\to$F: si risale con pred da F: F$\to$D$\to$B$\to$A, e Print-Path lo stampa al dritto: A -> B -> D -> F};
\end{tikzpicture}
\end{center}

\needspace{7cm}
\textbf{La DFS sullo stesso grafo}, con i timestamp $d/f$ scritti accanto a ogni vertice.
Si parte da A, si va al primo vicino bianco (B), poi al primo vicino bianco di B (D), e
così via finché si può; da F non c'è più niente di bianco e si torna indietro: F finisce
per primo (6), poi E, C, D, B e infine A (11).
\begin{center}
\begin{tikzpicture}
  \node[vert] (A) at (0,0) {A};
  \node[vert] (B) at (1.8,1.1) {B};
  \node[vert] (C) at (1.8,-1.1) {C};
  \node[vert] (D) at (3.6,1.1) {D};
  \node[vert] (E) at (3.6,-1.1) {E};
  \node[vert] (F) at (5.4,0) {F};
  \draw[arco,black!25!sfondo] (A)--(C) (D)--(F);
  \draw[albArco,-{Stealth[length=3mm]}] (A)--(B);
  \draw[albArco,-{Stealth[length=3mm]}] (B)--(D);
  \draw[albArco,-{Stealth[length=3mm]}] (D)--(C);
  \draw[albArco,-{Stealth[length=3mm]}] (C)--(E);
  \draw[albArco,-{Stealth[length=3mm]}] (E)--(F);
  \node[font=\footnotesize\ttfamily,left=1mm of A] {0/11};
  \node[font=\footnotesize\ttfamily,above=1mm of B] {1/10};
  \node[font=\footnotesize\ttfamily,above=1mm of D] {2/9};
  \node[font=\footnotesize\ttfamily,below=1mm of C] {3/8};
  \node[font=\footnotesize\ttfamily,below=1mm of E] {4/7};
  \node[font=\footnotesize\ttfamily,right=1mm of F] {5/6};
  \node[etichetta,anchor=west,text width=5.6cm] at (7,0.3) {In verde gli archi dell'albero DFS (A$\to$B$\to$D$\to$C$\to$E$\to$F): la DFS è andata \q{a fondo} lungo un solo ramo.};
  \node[etichetta,anchor=west,text width=5.6cm] at (7,-1.0) {A--C e D--F sono archi all'indietro: arrivano a vertici ancora grigi.};
\end{tikzpicture}
\end{center}
Confronta con la BFS: lì C era a distanza 1 da A; qui la DFS ci arriva dopo tre archi
(A, B, D, C). Gli intervalli sono annidati: $[3, 8]$ di C sta dentro $[2, 9]$ di D, che sta
dentro $[1, 10]$ di B, come le chiamate ricorsive.

\needspace{6cm}
\textbf{La pila delle chiamate} della DFS nel momento più profondo, quando si visita F
(Modulo 8): sei chiamate aperte una dentro l'altra, tutte con il vertice grigio.
\begin{center}
\begin{tikzpicture}
  \foreach \v [count=\i from 0] in {A,B,D,C,E,F}
    \node[cella,minimum width=3.4cm,fill=stackcol] at (0,0.62*\i) {dfsVisit(\v)};
  \node[etichetta,anchor=west] at (1.9,3.1) {$\leftarrow$ cima: F non ha vicini bianchi, finisce ($f=6$)};
  \node[etichetta,anchor=west] at (1.9,0) {$\leftarrow$ prima chiamata, dal ciclo di DFS(G)};
\end{tikzpicture}
\end{center}

% ---------------------------------------------------------------------
\section{Pattern di Codice di Laboratorio}
% ---------------------------------------------------------------------

Due classi template. \cc|GrafoMatrice<T>| segue il codice del professore: un array di
chiavi, una matrice \cc|bool**| e una ricerca lineare dalla chiave all'indice.
\cc|Grafo<T>| usa le liste di adiacenza e implementa BFS con Print-Path e DFS con i
timestamp esattamente come gli pseudocodici delle slide (colori compresi). Il \cc|main|
costruisce il grafo A--F in entrambe e lo visita.

\codice{codice/m14/pattern.cpp}
\uscita{codice/m14/output/pattern.txt}

\textbf{Cosa osservare:}
\begin{itemize}
  \item la matrice stampata è simmetrica, e ogni lista ha tanti elementi quanti sono gli 1
        della riga corrispondente;
  \item i valori stampati coincidono con le figure della sezione precedente: la BFS dà le
        distanze 0, 1, 1, 2, 2, 3; la DFS i tempi da 0 a 11;
  \item il cammino minimo verso F passa per B e non per C: dipende solo dall'ordine dei
        vicini nella lista di A (B viene prima). Entrambi i cammini hanno lunghezza 3;
  \item nella DFS la visita continua finché c'è qualcosa di bianco: se il grafo non fosse
        connesso, il ciclo di \cc|DFS()| ripartirebbe da un altro vertice;
  \item il distruttore libera prima i nodi di ogni lista, poi gli array: ogni \cc|new| ha
        il suo \cc|delete|. Copia e assegnazione sono vietate con \cc|= delete| (Regola dei
        Tre, Modulo 4).
\end{itemize}

% ---------------------------------------------------------------------
\section{Tips \& Trappole da Esame}
% ---------------------------------------------------------------------

\begin{trappola}[: dimenticare i vertici visitati]
Una BFS o DFS senza colori (o \cc|visitato[]|) su un grafo con un ciclo non termina: la
coda si riempie all'infinito o la ricorsione esaurisce lo stack. È l'errore più grave
sui grafi.
\end{trappola}

\begin{trappola}[: marcare in BFS quando si estrae, invece che quando si inserisce]
Il vertice va marcato (grigio) \textbf{nel momento in cui entra in coda}. Se lo marchi solo
quando esce, nel frattempo un altro vicino può rimetterlo in coda: viene visitato due
volte e le distanze possono uscire sbagliate.
\end{trappola}

\begin{trappola}[: grafo non orientato con un solo inserimento]
In un grafo non orientato l'arco $(u,v)$ va scritto \textbf{due volte}: in
\cc|adj[u][v]| e \cc|adj[v][u]|, oppure nella lista di $u$ e in quella di $v$. Altrimenti
hai costruito un grafo orientato, e la visita non trova metà dei cammini.
\end{trappola}

\begin{trappola}[: chiave e indice]
Il vertice \q{Catania} non è il numero 2: è \textbf{in posizione} 2. Le funzioni pubbliche
ricevono chiavi, all'interno si lavora con indici. Controlla sempre che \cc|indice()| non
restituisca $-1$ prima di usarlo per indicizzare un array.
\end{trappola}

\begin{trappola}[: una funzione che stampa dentro un \texttt{cout}]
\cc|cout << "Gruppi: " << g.componenti();| con una \cc|componenti()| che stampa a sua volta
produce un output mescolato: in pratica la funzione viene chiamata prima che la riga sia
scritta. Chiama la funzione prima, salva il risultato in una variabile e stampalo dopo.
È successo davvero scrivendo il laboratorio di questo modulo.
\end{trappola}

\begin{trappola}[: BFS da una sola sorgente su un grafo non connesso]
Una BFS da $s$ raggiunge solo la componente di $s$. Per visitare tutto il grafo serve il
ciclo esterno \q{per ogni vertice ancora bianco, riparti}, come in DFS($G$).
\end{trappola}

\begin{consiglio}[: lo schema fisso del distruttore]
Per le liste di adiacenza: \cc|for| su ogni vertice, \cc|while| che stacca e cancella i
nodi della sua lista, poi \cc|delete[]| di tutti gli array. Per la matrice: \cc|delete[]|
di ogni riga, poi \cc|delete[]| dell'array di righe. Al contrario dell'ordine di creazione.
\end{consiglio}

\begin{consiglio}[: simulare a mano prima di scrivere]
All'esame (e all'orale) ti chiederanno di eseguire a mano una BFS o una DFS. Usa la
tabella del passo per passo di questo modulo: colonna \q{estraggo}, colonna
\q{scopro}, colonna \q{coda}. Per la DFS, scrivi a matita $d/f$ accanto ai vertici.
\end{consiglio}

% ---------------------------------------------------------------------
\section{Quiz a Crocette (Prova Scritta)}
% ---------------------------------------------------------------------

\begin{quiz}{Un grafo non orientato ha 5 archi e viene rappresentato con liste di adiacenza. Quanti nodi contengono, in totale, tutte le liste?}
  \item 5
  \item 10
  \item 25
  \item dipende dal numero di vertici
\end{quiz}

\begin{quiz}{Quanto spazio occupa la matrice di adiacenza di un grafo con $|V|$ vertici e $|E|$ archi?}
  \item $O(|E|)$
  \item $O(|V| + |E|)$
  \item $O(|V|^2)$
  \item $O(|V| \cdot |E|)$
\end{quiz}

\begin{quiz}{Quale struttura dati di supporto usa la visita in ampiezza (BFS)?}
  \item una pila
  \item una coda
  \item un albero binario di ricerca
  \item nessuna
\end{quiz}

\begin{quiz}{In quale rappresentazione si verifica in tempo costante se esiste l'arco $(u,v)$?}
  \item liste di adiacenza
  \item matrice di adiacenza
  \item in entrambe
  \item in nessuna delle due
\end{quiz}

\begin{quiz}{Nel grafo del pattern (A: B C; B: A D; C: A D E; D: B C F; E: C F; F: D E), in che ordine la BFS da A estrae i vertici dalla coda?}
  \item A B D C E F
  \item A B C D E F
  \item A C E F D B
  \item A B C F D E
\end{quiz}

\begin{quiz}{Perché durante una visita i vertici vengono colorati (o marcati come visitati)?}
  \item per stampare il grafo a colori
  \item per evitare di visitare più volte lo stesso vertice e di girare all'infinito nei cicli
  \item per ordinare i vertici
  \item perché la coda accetta solo vertici colorati
\end{quiz}

\begin{quiz}{La matrice di adiacenza di un grafo non orientato è sempre:}
  \item diagonale
  \item simmetrica
  \item piena di 1
  \item triangolare inferiore
\end{quiz}

\begin{quiz}{Cosa calcola la BFS da $s$ su un grafo non pesato?}
  \item la distanza minima (in numero di archi) da $s$ a ogni vertice raggiungibile
  \item l'ordinamento topologico
  \item il ciclo più lungo
  \item il vertice di grado massimo
\end{quiz}

\begin{quiz}{Su quali grafi esiste un ordinamento topologico?}
  \item su tutti i grafi non orientati
  \item sui grafi orientati aciclici (DAG)
  \item solo sugli alberi binari
  \item sui grafi completi
\end{quiz}

\begin{quiz}{Qual è il costo di una DFS su un grafo rappresentato con liste di adiacenza?}
  \item $O(|V|)$
  \item $O(|E|^2)$
  \item $O(|V| + |E|)$
  \item $O(|V|^2)$
\end{quiz}

\begin{soluzione}[: risposte commentate]
\begin{enumerate}[label=\textbf{\arabic*.}]
  \item \textbf{B}. Ogni arco compare in due liste: $2 \times 5 = 10$.
  \item \textbf{C}. Sempre $|V| \times |V|$ celle, qualunque sia il numero di archi.
  \item \textbf{B}. La DFS usa invece una pila (o la ricorsione).
  \item \textbf{B}. Basta leggere \cc|adj[u][v]|; con le liste bisogna scorrere quella di $u$.
  \item \textbf{B}. Vedi la tabella del passo per passo. La A è l'ordine della DFS.
  \item \textbf{B}.
  \item \textbf{B}. L'arco $(u,v)$ si scrive sia in \cc|adj[u][v]| sia in \cc|adj[v][u]|.
  \item \textbf{A}.
  \item \textbf{B}. Un ciclo renderebbe impossibile rispettare tutti i vincoli.
  \item \textbf{C}. Ogni vertice visitato una volta, ogni lista scorsa una volta.
\end{enumerate}
\end{soluzione}

% ---------------------------------------------------------------------
\section{Laboratorio a Step (Simulazione d'Esame)}
% ---------------------------------------------------------------------

\begin{traccia}[: rete di amicizie e piano di studi]
Implementare una classe template \cc|Grafo<T>| con liste di adiacenza, orientato o non
orientato a scelta del costruttore, senza usare la STL. Usarla per: (a) una rete di
amicizie (grafo non orientato di nomi): numero di amici di una persona e persona più
popolare, distanza tra due persone, catena di amicizie più corta, gruppi di persone
collegate; (b) un piano di studi (grafo orientato di esami, arco da A a B se A è
propedeutico a B): un ordine valido in cui sostenere tutti gli esami.
\end{traccia}

\begin{step}{la classe Grafo}
Scrivi \cc|Grafo<T>| con l'array delle chiavi, l'array \cc|Arco**| delle teste di lista,
il flag \cc|orientato|, \cc|indice|, \cc|aggiungiVertice|, \cc|aggiungiArco| (che restituisce
\cc|false| se un vertice non esiste), \cc|stampa| e distruttore. Inserisci i vicini in
testa.\\[0.3em]
\textbf{Concetto nuovo:} array di liste concatenate.
\end{step}

\begin{step}{gradi}
Scrivi \cc|grado(k)| (lunghezza della lista di $k$) e \cc|gradoMassimo()|, la persona con
più amici.\\[0.3em]
\textbf{Concetto nuovo:} leggere le informazioni dalle liste di adiacenza.
\end{step}

\begin{step}{distanza con la BFS}
Scrivi una \cc|BFS| che riempie \cc|dist[]| e \cc|pred[]| usando una \cc|Coda<int>| su nodi
dinamici, e \cc|distanza(a, b)|, che restituisce $-1$ se $b$ non è raggiungibile.\\[0.3em]
\textbf{Concetto nuovo:} visita in ampiezza e distanze minime.
\end{step}

\begin{step}{catena più corta}
Scrivi \cc|cammino(a, b)|, che dopo la BFS stampa la catena di amicizie da $a$ a $b$ e la
sua lunghezza, usando una funzione ricorsiva privata sul modello di Print-Path.\\[0.3em]
\textbf{Concetto nuovo:} ricostruire un cammino dai predecessori.
\end{step}

\begin{step}{gruppi con la DFS}
Scrivi \cc|componenti()|, che stampa i gruppi di persone collegate (componenti connesse)
con una DFS ricorsiva e ne restituisce il numero.\\[0.3em]
\textbf{Concetto nuovo:} visita in profondità ripetuta da ogni vertice non visitato.
\end{step}

\begin{step}{ordinamento topologico}
Costruisci il grafo orientato degli esami e scrivi \cc|ordinamentoTopologico()|: una DFS
che, quando un vertice finisce, lo inserisce in testa a una lista concatenata; poi stampa
e libera la lista.\\[0.3em]
\textbf{Concetto nuovo:} usare il tempo di fine della DFS.
\end{step}

\subsection*{Soluzione completa}
\codice{codice/m14/lab.cpp}
\uscita{codice/m14/output/lab.txt}

\begin{consiglio}[: cosa notare nella soluzione]
\begin{itemize}
  \item Con l'inserimento in testa le liste risultano \textbf{rovesciate} rispetto
        all'ordine degli archi (Anna: Carla Bruno). È per questo che la catena verso
        Fabio passa per Carla e non per Bruno: entrambe sono lunghe 4, la BFS sceglie la
        prima che trova.
  \item Da Bruno a Carla la catena è lunga 2; ce ne sono due (via Anna o via Dario) e la
        BFS ne restituisce una sola.
  \item \cc|visitato[]|, \cc|dist[]| e \cc|pred[]| sono membri della classe, allocati una
        volta nel costruttore e riusati a ogni visita: per questo ogni visita li
        reinizializza all'inizio.
  \item Nell'ordinamento topologico ogni esame compare \textbf{dopo} i suoi propedeutici
        (Programmazione 1 prima di Programmazione 2, che viene prima di Basi di Dati e di
        Algoritmi), anche se l'ordine di inserimento degli esami era quasi al contrario.
        L'ordine non è unico: ne esistono molti validi.
  \item Per verificare a mano: ogni arco della stampa del grafo orientato deve andare da
        un esame a uno che compare più in basso nell'elenco finale.
\end{itemize}
\end{consiglio}

% ---------------------------------------------------------------------
\section{Domande e Risposte Modello per l'Orale}
% ---------------------------------------------------------------------

\begin{orale}[D1. Cos'è un grafo? Che differenza c'è tra orientato e non orientato?]
Un grafo $G = (V, E)$ è formato da un insieme di vertici $V$ e da un insieme di archi $E$,
ciascuno dei quali è una coppia di vertici. In un grafo orientato la coppia è ordinata:
l'arco $(u, v)$ esce da $u$ ed entra in $v$, e non implica l'arco $(v, u)$. In un grafo non
orientato le coppie non sono ordinate, quindi l'adiacenza è simmetrica, e nelle slide non
sono ammessi cappi. Diversamente da un albero, un grafo può contenere cicli e non ha una
radice; un albero è un caso particolare di grafo non orientato, connesso e aciclico.
\end{orale}

\begin{orale}[D2. Come si rappresenta un grafo in memoria? Quando conviene ciascuna rappresentazione?]
Con la matrice di adiacenza, una matrice $|V| \times |V|$ in cui l'elemento $(i,j)$ vale 1
se c'è l'arco da $i$ a $j$: occupa $O(|V|^2)$ e permette di sapere in tempo costante se
due vertici sono adiacenti; per un grafo non orientato è simmetrica. Oppure con le liste
di adiacenza, un array di $|V|$ liste in cui la lista del vertice $u$ contiene i suoi
adiacenti: occupa $O(|V| + |E|)$, la somma delle lunghezze delle liste è $|E|$ per un
grafo orientato e $2|E|$ per uno non orientato. Le liste convengono per grafi sparsi e per
gli algoritmi che scorrono i vicini, come le visite; la matrice per grafi densi o quando
servono molte verifiche di adiacenza.
\end{orale}

\begin{orale}[D3. Descriva la visita in ampiezza.]
La BFS parte da una sorgente $s$ e scopre i vertici in ordine di distanza da $s$. Tutti i
vertici sono inizialmente bianchi; $s$ diventa grigio con distanza 0 e viene messo in una
coda. Finché la coda non è vuota si estrae un vertice $u$ e, per ogni suo adiacente
bianco $v$, lo si colora di grigio, si pone $d[v] = d[u] + 1$ e $\text{pred}[v] = u$ e lo
si accoda; finiti gli adiacenti, $u$ diventa nero. La coda contiene sempre i vertici
grigi. Alla fine $d[v]$ è la lunghezza del cammino minimo da $s$ a $v$ e i predecessori
formano il breadth-first tree, da cui la procedura ricorsiva Print-Path ricostruisce il
cammino. Con le liste di adiacenza il costo è $O(|V| + |E|)$.
\end{orale}

\begin{orale}[D4. Descriva la visita in profondità e un suo uso.]
La DFS esplora a partire dal vertice scoperto più di recente che ha ancora archi non
esplorati; quando non ce ne sono più torna indietro (backtracking). Si implementa con una
procedura ricorsiva DFS-Visit, chiamata da DFS($G$) su ogni vertice ancora bianco, per cui
il risultato è una foresta. Ogni vertice riceve due timestamp: $d[v]$ quando viene
scoperto e $f[v]$ quando ha finito di esaminare la sua lista. Il costo è
$\Theta(|V| + |E|)$. Un uso è l'ordinamento topologico di un grafo orientato aciclico:
si esegue la DFS e ogni vertice, quando finisce, viene inserito in testa a una lista; si
ottengono i vertici in ordine decrescente di $f$, e ogni arco va da un vertice a uno
successivo. Altri usi sono la ricerca di cicli e le componenti connesse.
\end{orale}

\begin{orale}[D5. Perché nelle visite di un grafo servono i colori, mentre negli alberi no?]
Perché un grafo può contenere cicli e più cammini verso lo stesso vertice. Senza tenere
traccia dei vertici già scoperti, la visita ripasserebbe più volte dagli stessi vertici e,
in presenza di un ciclo, non terminerebbe. In un albero ogni nodo ha un solo padre e
scendendo non si torna mai a un nodo già visto, quindi il problema non si pone. I colori
garantiscono anche che ogni vertice venga esaminato una sola volta, da cui il costo
lineare $O(|V| + |E|)$.
\end{orale}

% ---------------------------------------------------------------------
\section{Esercizi al PC}
% ---------------------------------------------------------------------

\begin{esercizio}{Matrice di adiacenza}{facile}
Con una matrice statica \cc|bool m[5][5]| costruisci il grafo non orientato con archi 0--1,
0--2, 1--2, 2--3, 3--4. Stampa la matrice con il grado di ogni vertice, il numero di
archi, se 0 e 2 (e 0 e 4) sono adiacenti, e se il grafo è completo. Poi aggiungi tutti gli
archi mancanti e verifica che diventi completo con $5 \cdot 4 / 2 = 10$ archi.
\end{esercizio}

\begin{esercizio}{Liste di adiacenza a mano}{facile}
Con un array di 4 teste \cc|Arco*| costruisci lo stesso insieme di 4 archi una volta come
grafo non orientato e una volta come orientato. Stampa le liste nella forma
\texttt{0: -> 2 -> 1 -> null} e verifica che la somma delle lunghezze sia $2|E|$ nel primo
caso e $|E|$ nel secondo. Libera tutta la memoria.
\end{esercizio}

\begin{esercizio}{Da matrice a liste e ritorno}{facile}
Data la matrice di un grafo orientato, costruisci le liste di adiacenza in modo che i
vicini risultino in ordine crescente (pur inserendo in testa). Poi ricostruisci la
matrice dalle liste e controlla che sia uguale all'originale.
\end{esercizio}

\begin{esercizio}{Gradi e grafo trasposto}{medio}
In un grafo orientato con liste di adiacenza calcola, con un solo passaggio su tutte le
liste, il grado entrante e uscente di ogni vertice. Poi costruisci il grafo trasposto
$G^T$, in cui ogni arco $(u,v)$ diventa $(v,u)$.
\end{esercizio}

\begin{esercizio}{Amici di amici}{medio}
In una rete di 9 persone (una isolata), stampa con una BFS da \q{Tu} le persone raggruppate
per distanza (\texttt{distanza 1: Ada Bea}, \ldots), e alla fine quelle non raggiungibili.
Suggerimento: i vertici escono dalla coda in ordine di distanza; vai a capo quando la
distanza cambia.
\end{esercizio}

\begin{esercizio}{Esiste un cammino?}{medio}
Scrivi due funzioni che dicono se esiste un cammino da $s$ a $t$: una con la BFS e una con
una DFS \textbf{iterativa} che usa una pila esplicita al posto della ricorsione. Stampa
l'ordine di visita della DFS iterativa e provale su un grafo formato da due pezzi
separati.
\end{esercizio}

\begin{esercizio}{Componenti connesse}{medio}
Con una DFS ricorsiva che restituisce quanti vertici ha marcato, assegna a ogni vertice il
numero della sua componente (\cc|comp[u]|). Stampa ogni componente con la sua dimensione,
dì se il grafo è connesso e usa \cc|comp[]| per rispondere in $O(1)$ a domande come
\q{4 e 8 sono collegati?}.
\end{esercizio}

\begin{esercizio}{C'è un ciclo?}{difficile}
In un grafo orientato, cerca un ciclo con una DFS a tre colori: un arco verso un vertice
\textbf{grigio} (antenato ancora aperto) chiude un ciclo; un arco verso un vertice
\textbf{nero} no. Provala su un \q{rombo} senza cicli (0$\to$1, 0$\to$2, 1$\to$3, 2$\to$3,
3$\to$4) e su un grafo con il ciclo 1$\to$2$\to$3$\to$1. Perché con due soli colori il
rombo verrebbe giudicato ciclico?
\end{esercizio}

\begin{esercizio}{Uscire dal labirinto}{difficile}
Un labirinto è una griglia di caratteri (\texttt{\#} muro, \texttt{.} libero). Ogni cella
libera è un vertice, adiacente alle celle libere sopra, sotto, a destra e a sinistra: il
grafo non va costruito, i vicini si calcolano al volo. Con una BFS trova il percorso più
breve dall'angolo in alto a sinistra a quello in basso a destra e disegnalo con delle
\texttt{o}. Suggerimento: numera la cella $(r, c)$ come $r \cdot C + c$.
\end{esercizio}

\begin{esercizio}{Grafo bipartito}{difficile}
Un grafo è \textbf{bipartito} se si possono colorare i vertici con due colori in modo che
ogni arco unisca vertici di colore diverso (esempio: studenti e corsi, con un arco
quando uno studente segue un corso). Verificalo con una BFS che dà a ogni vertice scoperto
il colore opposto a quello di chi lo ha scoperto, e si ferma se trova un arco con estremi
dello stesso colore. Provala su un esagono (sì) e su un triangolo (no).
\end{esercizio}

\subsection*{Soluzioni}
Le soluzioni commentate, con il codice e l'output reale, sono nel PDF separato
\textbf{Soluzioni} (\texttt{soluzioni-light.pdf} / \texttt{soluzioni-dark.pdf}),
allo stesso numero di esercizio. Il codice è anche in \texttt{codice/m14/}.
